%% kapa: exam=1 round=21 year=2010 period=2 form=A subject=appraisal_law subject_ko=부동산관계법규 seq=25 number=65
%% source: Q-Net 감정평가사 기출문제 (cst00309) artlSeq=5024181 file=extracted/1008/5024181/2010년 제21회 감정평가사 1차시험 기출문제/2교시.hwp
%% conversion: kordoc 4.24.51 -> tools/gen1.py (auto); review=proofread-gpt-sol; flags=none
\begin{question}{65}
국유재산법령상 국유재산에 관한 설명으로 옳지 않은 것은?
\begin{choices}
  \item 행정재산은「민법」제245조(점유로 인한 부동산소유권의 취득기간)에도 불구하고 시효취득의 대상이 되지 아니한다.
  \item 판결에 따라 취득하는 경우를 제외하고는 사권이 설정된 재산은 그 사권이 소멸된 후가 아니면 국유재산으로 취득하지 못한다.
  \item 행정재산에 대한 종전의 사용허가를 갱신받으려는 자는 허가기간이 끝나기 1개월 전에 관리청에 신청하여야 한다.
  \item 일반재산 중 조림을 목적으로 하는 토지와 그 정착물의 경우에는 대부기간을 20년 이내로 한다.
  \item 관리청은 행정재산의 사용허가를 받은 자가 관리청의 승인 없이 사용허가를 받은 재산의 원래 상태를 변경하는 경우에는 그 허가를 취소하거나 철회할 수 있다.
\end{choices}
\end{question}
\answer{4}
